\chapter*{阅读与使用}
\addcontentsline{toc}{chapter}{阅读与使用}

本书以理论建模、数学推导和图像分析为主线。阅读一章的目标，是能说明一个模型适用的条件，能从图中建立方程，并能检查结论，而不只是记住最后一个公式。题图给出对象、约束和符号；解答图进一步标出分量、微元、过程或极限。阅读时应把文字、公式和图来回对照，特别留意正方向、参考系、积分区域以及状态与过程的区别。

\section*{先建立基础，再进入选读}
数学工具章可按当前专题的需要反复查阅。学习运动学与力学时，先熟悉位置、速度、受力图和守恒条件，再使用微分、积分表达连续变化。接着学习振动、引力、波和流体，把前面的方法用于新的约束与边界。热学中先分清状态量、过程量和功热符号，再学习概率、分布与涨落；电磁学中先建立场、势、电路与感应的基本关系，再进入传播和材料响应。光学与近代物理既延续这些方法，也会引入新的模型假设。

标有“选读”的部分用于扩展理论视野。第一遍可以先掌握相应章节的基本模型和常规例题，再回头阅读多自由度振动、信息论或场论等较深内容。高阶选读需要更熟练的微积分、线性代数和相关物理基础。高级定义出现得多，并不意味着基础已经掌握；能否清楚处理一个新的简单情形，往往比能否复述复杂公式更能检验理解。

\section*{用例题训练独立推导}
先遮住解答，把题设改写成自己的图：哪些量已知，哪些待求，哪些边界保持固定，是否允许接触、脱离、滑动或穿越。再选研究系统与参考系，写出控制关系，并说明采用的近似。做到这一步后再阅读解答，比较的是模型选择与推导次序，不能只对照末行数值。

完成计算后，至少检查量纲、符号与一个能直接理解的极限。求到代数根时还要检查定义域与原方程；求到平衡点时还要区别存在性与稳定性；使用无穷过程或小量展开时，应说明极限的取法与误差的阶。发现题图与算式不一致时，先回到定义核对，而不是为了迎合一个熟悉答案而临时改变正方向。

% CompletenessTail20261024 guide-and-notation
\par\Needspace{.72\textheight}
\section*{把理论理解接到实际训练}
例题的图和连续曲线表示模型中的理想关系，实际测量还包含仪器分辨率、随机起伏、系统偏差和模型失配。理论学习应配合独立作答、限时表达与实验数据处理。涉及实验的内容，需要实际进行器材读数、方案搭建、重复测量和异常分析；只阅读原理图不能代替操作训练。

教师可根据学生基础安排学习次序与选读深度。每次学习应留下可检查的产物：一张标清条件的图、一段有依据的推导，以及对结果有效范围的说明。这样才能把读懂例题逐步变成在陌生问题中建立模型的能力。

% Completeness20261024: put the notation beside the reading guide.
\section*{记号与单位约定}
\phantomsection
\addcontentsline{toc}{section}{记号与单位约定}
除题目另有说明外，本书采用国际单位制。数值计算所用物理常数和近似值以题目或所在节的说明为准。

矢量通常以粗体表示，部分图示与推导使用箭头记号；两者均表示具有大小和方向的量。
微分记号写作 $\dd$，虚数单位写作 $\ii$。
同一字母在不同专题中可能表示不同物理量，应结合本节定义及单位辨认；
例如 $I$ 可表示电流、转动惯量或光强，$L$ 可表示长度、角动量大小或电感。

除专门说明角度单位的题目外，微积分与小角展开中的角度以弧度计。
对数的底、正方向及参考系均以局部定义为准。各章例题、图和公式均按“章号.顺序号”编号。

